\subsubsection{Spin and trajectory detection results}

The ball detection pipeline was applied to eight clips across two release sessions. The releases of clips 9-12 and 13-16 were done consecutively with a non-homogeneous ball. For each clip, the lane polygon was reused for each session, across all clips within it. The \emph{Circle Hough Transform} consistently detected the ball across all frames after preprocessing, with the weighted directed graph resolving ambiguous candidates into a coherent single trajectory. Cubic spline interpolation and Gaussian smoothing ($\sigma = 5$) filled the sparse detection gaps (as seen in Figure \ref {fig:final_tracking}), and the final radius estimates followed the expected monotonic decrease as the ball travels away from the camera.

\begin{figure}[H]
    \centering
    \includegraphics[height=0.25\textwidth]{ball_tracking/final_tracking_1.png}
    \includegraphics[height=0.25\textwidth]{ball_tracking/final_tracking_2.png}
    \includegraphics[height=0.25\textwidth]{ball_tracking/final_tracking_3.png}
    \includegraphics[height=0.25\textwidth]{ball_tracking/final_tracking_4.png}
    \caption{Final ball tracking visualization with and witouth occluded ball}
    \label{fig:final_tracking}
\end{figure}

The spin detection workflow was executed on all eight clips, yielding per-frame estimates of the rotation axis and angular velocity as shown in Table \ref{tab:spin_results}. The angular velocity $\omega$ is computed per consecutive frame pair via the Kabsch algorithm and averaged across the clip. The rev rate is derived as:
\begin{equation}
	RPM = \dfrac{\omega \cdot fps }{2\pi}\cdot 60
\end{equation}


\begin{table}[htbp]
\centering
    \caption{Spin detection results across two recording sessions.}
    \label{tab:spin_results}
    \small
    \setlength{\tabcolsep}{4pt}
    \begin{tabular}{ccccccc}
    \toprule
    \textbf{Session} & \textbf{Clip} & \textbf{Frames} & \textbf{FPS} & \textbf{Avg. $\omega$ (rad/f)} & \textbf{RPM} & \textbf{Mean Rotation Axis} \\
    \midrule
    \multirow{4}{*}{A} & clip\_9  & 98  & 30 & 0.5470 & 156.7 & [0.46, 0.07, -0.88] \\
                    & clip\_10 & 125 & 30 & 0.5025 & 144.0 & [0.00, 0.99, 0.12] \\
                    & clip\_11 & 207 & 30 & 0.5517 & 158.1 & [-0.03, 0.12, -0.99] \\
                    & clip\_12 & 109 & 30 & 0.5183 & 148.5 & [0.87, 0.26, -0.41] \\
    \midrule
    \multirow{4}{*}{B} & clip\_13 & 196 & 30 & 0.2970 & 85.1 & [0.18, -0.52, 0.83] \\
                    & clip\_14 & 118 & 30 & 0.3628 & 103.9 & [0.23, 0.81, -0.54] \\
                    & clip\_15 & 104 & 30 & 0.4095 & 117.3 & [-0.06, -0.60, -0.79] \\
                    & clip\_16 & 138 & 30 & 0.4092 & 117.2 & [0.12, -0.97, 0.20] \\
    \bottomrule
    \end{tabular}%
\end{table}

The two sessions exhibit clearly distinct rev rate profiles. Session A yields a mean of \emph{151.8 RPM} with a standard deviation of \emph{6.7 RPM}, while Session B yields a mean of \emph{105.9 RPM} with a standard deviation of \emph{15.2 RPM}. The standard deviations are 4.4\% and 14.4\% of their respective means. Despite a slightly higher variance in Session B, there remains a strong overall in-session consistency across clips of varying length. This repeatability suggests the pipeline is responding to a stable physical characteristic of each release rather than to measurement noise.

The separation between sessions is consistent across all clips with no overlap: the lowest estimate in Session A (144.0 RPM, Clip 10) exceeds the highest estimate in Session B (117.3 RPM, Clip 15) by 26.7 RPM. This clean separation further supports the interpretation that the two sessions correspond to physically distinct release styles or bowlers, with Session A exhibiting a meaningfully higher rev rate.

Furthermore, the mean 3D rotation axis provides a robust mathematical fingerprint of the bowler's specific axis tilt and axis rotation at release. Because the postprocessing pipeline incorporates a vector unflipping step to resolve sign ambiguities inherent to 3D motion estimation, the computed mean axis remains highly stable throughout the ball's trajectory. The clear variations in the rotation axis components between clips further corroborate that the two sessions captured physically distinct release characteristics.

\begin{figure}[H]
    \centering
    \includegraphics[height=0.25\textwidth]{ball_tracking/spin_visualization_1.png}
    \includegraphics[height=0.25\textwidth]{ball_tracking/spin_visualization_2.png}
    \includegraphics[height=0.25\textwidth]{ball_tracking/spin_visualization_3.png}
    \includegraphics[height=0.25\textwidth]{ball_tracking/spin_visualization_4.png}
    \caption{Final visualization of ball's spin (yellow and blue) and rotation axis (red). Note that the rendered ball texture does not perfectly align the axis as the rendering starts from a different orientation than the axis}
    \label{fig:spin_visualization}
\end{figure}